%=====================================================================
\chapter{From Virtual Machines to Containers}\label{ch:containers}
%=====================================================================
\locator{3}{Part III --- Containers and Lightweight Virtualization}

\begin{outcomes}
By the end of this chapter you will be able to:
\begin{tight}
  \item \textbf{Explain} where a container cuts the stack, and \textbf{contrast}
        it with where a hypervisor cuts it.
  \item \textbf{Name} the Linux namespaces and \textbf{state} what each one
        isolates.
  \item \textbf{Use} control groups to limit a process tree's CPU and memory,
        and \textbf{explain} what happens when a limit is reached.
  \item \textbf{Describe} how an overlay filesystem gives every container a
        private root from shared read-only layers.
  \item \textbf{Compute} the density and start-up advantage of containers over
        virtual machines, and \textbf{state} precisely what is given up for it.
\end{tight}
\end{outcomes}

\begin{prereq}
\chref{architectures} (the taxonomy figure, which placed containers on the
second row) and \chref{cpumem} (overcommit and the working set). From Operating
Systems you need processes, the process tree, file descriptors and mount
points. Nothing in Part~II's hardware material is required --- which is itself
the point of this chapter.
\end{prereq}

\begin{keyterms}
OS-level virtualization $\bullet$ container $\bullet$ namespace $\bullet$ mount
namespace $\bullet$ PID namespace $\bullet$ network namespace $\bullet$ user
namespace $\bullet$ control group $\bullet$ cgroup v2 $\bullet$ cgroup
controller $\bullet$ union filesystem $\bullet$ overlayfs $\bullet$ lower dir
$\bullet$ upper dir $\bullet$ copy-up $\bullet$ capability $\bullet$ seccomp
$\bullet$ chroot $\bullet$ rootless container $\bullet$ shared kernel
\end{keyterms}

\section{A Different Place to Cut}

Every chapter so far has virtualized \emph{hardware}: a processor that is not a
processor, memory that is not where the guest thinks, a disk that is a file.
The guest runs its own kernel on that imaginary machine.

A container asks a cheaper question. If all the workloads are going to run on
Linux anyway, why buy forty copies of Linux? Leave one kernel in place, and
give each group of processes a private \emph{view} of it: its own filesystem
root, its own process numbers, its own network interfaces, its own users. The
processes are ordinary processes on the host kernel. They simply cannot see
anything they were not shown.

\dsfig{%
\begin{tikzpicture}[font=\scriptsize,
  bx/.style={rounded corners=2pt, line width=0.8pt, align=center,
             inner sep=2pt, minimum height=0.5cm, font=\scriptsize},
  hw/.style={bx, draw=neutral, fill=neutral!15, text=neutral!40!black,
             minimum width=6.4cm},
  hv/.style={bx, draw=primary, fill=primary!12, text=primary,
             minimum width=6.4cm},
  os/.style={bx, draw=purplehl, fill=purplehl!12, text=purplehl!70!black},
  ap/.style={bx, draw=secondary, fill=secondary!8, text=secondary!80!black},
  cut/.style={accent, line width=1.3pt, dashed},
  note/.style={font=\scriptsize\itshape, text=accent, align=center}
]
  % ---- virtual machines ----
  \node[font=\small\bfseries, text=primary] at (3.2,4.65) {Virtual machines};
  \node[hw] at (3.2,0)    {hardware};
  \node[hv] at (3.2,0.72) {hypervisor};
  \foreach \i/\x in {0/1.1, 1/3.2, 2/5.3}{
    \node[os, minimum width=1.9cm] at (\x,1.65) {guest kernel};
    \node[os, minimum width=1.9cm, fill=purplehl!6] at (\x,2.25)
         {libraries};
    \node[ap, minimum width=1.9cm] at (\x,2.85) {app};}
  \draw[cut] (-0.1,1.22) -- (6.5,1.22);
  \node[note] at (3.2,3.55) {three kernels, 200--500 MB\\of overhead each};

  % ---- containers ----
  \node[font=\small\bfseries, text=primary] at (11.6,4.65) {Containers};
  \node[hw] at (11.6,0)    {hardware};
  \node[os, minimum width=6.4cm] at (11.6,0.72) {\textbf{one} host kernel};
  \foreach \x in {9.5, 11.6, 13.7}{
    \node[os, minimum width=1.9cm, fill=purplehl!6] at (\x,1.65) {libraries};
    \node[ap, minimum width=1.9cm] at (\x,2.25) {app};}
  \draw[cut] (8.3,1.22) -- (14.9,1.22);
  \node[note] at (11.6,3.55) {one kernel, a few MB\\of overhead each};

  \node[font=\scriptsize\bfseries, text=accent, anchor=west] at (6.7,1.22)
       {the cut};
  \node[font=\scriptsize\bfseries, text=accent, anchor=east] at (8.2,1.22)
       {the cut};

  \node[rounded corners=3pt, draw=accent!60, fill=accent!5,
        text=accent!70!black, font=\scriptsize, align=center, inner sep=4pt,
        text width=12.6cm] at (7.4,-1.1)
       {Above the dashed line is what each workload gets privately; below it is
        what they share. Moving the line up buys density and spends
        isolation.};
\end{tikzpicture}}%
{The whole of Part III in one picture. Nothing else about containers is a new
idea --- only the height at which the stack is cut.}{where-the-cut-is}

\begin{definitionbox}[Container]
A \term{container} is a group of processes on a shared kernel, given a private
view of the system by \term{namespaces} and a bounded share of its resources by
\term{control groups}, usually running from a filesystem image assembled by a
\term{union filesystem}.

\smallskip
There is no ``container'' object in the Linux kernel. The kernel has
namespaces, cgroups, capabilities and mounts; a container is a convention about
using them together, and Docker's contribution was the convention and the
packaging, not the mechanism.
\end{definitionbox}

\section{Namespaces}

A namespace wraps one global system resource so that processes inside it see
their own instance of it. Eight exist, added over eighteen years.

\rowcolors{2}{white}{lightbg}
\begin{longtable}{@{}L{2.5cm}L{1.3cm}L{5.4cm}L{5.5cm}@{}}
\toprule
\rowcolor{primary!15}
\textbf{Namespace} & \textbf{Year} & \textbf{What it isolates} &
\textbf{Consequence inside} \\
\midrule
\endhead
\term{mount} & 2002 & The set of mount points & A private filesystem tree ---
this is the one that makes an image work \\
\term{UTS} & 2006 & Hostname and domain name & \texttt{hostname} returns
something of its own \\
\term{IPC} & 2006 & System V IPC, POSIX message queues & Shared memory
segments cannot be reached from outside \\
\term{PID} & 2008 & Process identifiers & The first process is PID 1, and the
host's processes are invisible \\
\term{network} & 2009 & Interfaces, routes, firewall rules, ports & Its own
\texttt{lo} and \texttt{eth0}; two containers may both bind port 80 \\
\term{user} & 2013 & User and group identifiers & \texttt{root} inside can map
to an unprivileged user outside --- the basis of rootless containers \\
\term{cgroup} & 2016 & The cgroup hierarchy root & A container cannot see where
it sits in the host's resource tree \\
\term{time} & 2020 & The boot and monotonic clocks & Used mainly for
checkpoint and restore \\
\bottomrule
\end{longtable}

\begin{conceptbox}[PID 1 Is Not a Detail]
Inside a PID namespace the first process is numbered 1, and the kernel treats
PID 1 specially everywhere: it must reap orphaned children, and signals with
default actions are \emph{not} delivered to it unless it has installed a
handler.

\smallskip
Two consequences bite in production. A container whose PID 1 is an ordinary
application does not reap orphans, so a process that forks accumulates zombies
until the PID limit is reached. And \texttt{docker stop} sends
\texttt{SIGTERM}, which an application that installed no handler simply
ignores --- the runtime then waits ten seconds and sends \texttt{SIGKILL},
which is why so many containers take exactly ten seconds to stop and lose
in-flight work.

\smallskip
Both are fixed by running a tiny init as PID 1, which is what
\texttt{docker run --init} does. Knowing \emph{why} is the difference between
applying the flag and diagnosing the symptom.
\end{conceptbox}

\begin{alertbox}[The user namespace is the one that changed the security story]
Before user namespaces, \texttt{root} inside a container was \texttt{root} on
the host: the same UID 0, with the same power over the shared kernel, and one
escape was a total compromise.

\smallskip
A user namespace maps identifiers. UID 0 inside can be UID 100000 outside, so a
process that believes it is root holds no privilege the host recognises. This
is what makes \term{rootless containers} possible --- an unprivileged user
running containers with no daemon and no setuid helper, which is Podman's
default and now an option in Docker.

\smallskip
It is not enabled by default in every runtime, and a container run as root
without it is still a process running as real root beside your other services.
\chref{security} returns to this.
\end{alertbox}

\section{Control Groups}

Namespaces decide what a container can \emph{see}. They say nothing about how
much it can \emph{use}: a process in a private namespace can still consume
every core and all the memory on the host.

\term{Control groups} supply the limits. Processes are arranged in a tree, and
each node carries controllers --- \texttt{cpu}, \texttt{memory}, \texttt{io},
\texttt{pids} --- that account for and cap consumption. Version~2, now the
default everywhere, replaced v1's separate per-controller hierarchies with a
single unified tree.

\begin{conceptbox}[What a Limit Does When It Is Reached]
The two important limits fail in completely different ways, and conflating them
is a common source of mystification.

\smallskip
\textbf{CPU is throttled.} A container limited to 0.5 of a core that asks for
more is simply scheduled less: it runs at half speed. Nothing fails. You see it
as latency, and in \texttt{cpu.stat} as a rising \texttt{throttled\_usec}.

\smallskip
\textbf{Memory is killed.} A container that exceeds its memory limit has a
process killed by the OOM killer --- usually the largest, usually the one you
cared about. There is no graceful degradation, because there is nothing to
degrade: the page is needed or it is not.

\smallskip
This is the same asymmetry as \chref{cpumem}: time is divisible and memory is
not. It is why the commonest container incident is an application restarting
every few hours with exit code 137 --- which is $128 + 9$, signal 9, the OOM
kill --- and why memory limits deserve far more thought than CPU limits.
\end{conceptbox}

\section{The Image, and the Union Filesystem}

A container needs a root filesystem. Giving each one a private copy would cost
hundreds of megabytes apiece and lose most of the density the chapter has been
claiming.

\term{Overlayfs} solves it by stacking directories. Several read-only
\term{lower} layers are presented as one tree, with a single writable
\term{upper} layer on top. A read finds the file in the topmost layer that has
it. A write triggers a \term{copy-up}: the file is copied into the upper layer
and modified there, leaving the shared layer untouched.

\dsfig{%
\begin{tikzpicture}[font=\scriptsize,
  lyr/.style={rounded corners=2pt, line width=0.8pt, align=left,
              minimum width=5.4cm, minimum height=0.56cm, inner sep=4pt,
              font=\scriptsize},
  ro/.style={lyr, draw=secondary, fill=secondary!8, text=secondary!80!black},
  rw/.style={lyr, draw=greenhl, fill=greenhl!12, text=greenhl!45!black},
  ct/.style={rounded corners=2pt, draw=purplehl, line width=0.9pt,
             fill=purplehl!8, text=purplehl!70!black, align=center,
             minimum width=1.9cm, minimum height=0.56cm, font=\scriptsize},
  note/.style={font=\scriptsize\itshape, text=neutral, anchor=west, align=left}
]
  \node[ro] at (0,0)    {\texttt{ubuntu:22.04} \quad 77 MB};
  \node[ro] at (0,0.66) {\texttt{+ python3} \quad 42 MB};
  \node[ro] at (0,1.32) {\texttt{+ our dependencies} \quad 18 MB};
  \node[ro] at (0,1.98) {\texttt{+ our application} \quad 2 MB};
  \node[rw] at (0,2.74) {writable layer --- \emph{per container}, empty at start};

  \node[note] at (3.0,0)    {shared by every container};
  \node[note] at (3.0,0.66) {built from this image};
  \node[note] at (3.0,1.32) {stored once on disk};
  \node[note] at (3.0,1.98) {and once in page cache};

  \foreach \i/\x in {1/7.6, 2/9.8, 3/12.0}{
    \node[ct] at (\x,2.74) {container \i};
    \draw[-{Stealth[length=1.6mm]}, purplehl!60, line width=0.7pt]
          (\x,2.42) -- (\x,2.1);}
  \node[rounded corners=3pt, draw=secondary!50, fill=secondary!4,
        minimum width=6.4cm, minimum height=1.9cm] at (9.8,1.0) {};
  \node[font=\scriptsize, text=secondary!80!black, align=center] at (9.8,1.0)
       {one copy of the four\\read-only layers,\\139 MB in total};

  \node[rounded corners=2pt, draw=accent, fill=accent!6, text=accent!70!black,
        font=\scriptsize, align=left, inner sep=4pt, text width=6.0cm]
       at (9.6,-1.25)
       {\textbf{Copy-up.} Container 2 edits
        \texttt{/etc/app.conf}: the file is copied into \emph{its} writable
        layer and changed there. Containers 1 and 3 still see the original.};
  \node[rounded corners=2pt, draw=neutral!55, fill=neutral!6,
        text=neutral!40!black, font=\scriptsize, align=left, inner sep=4pt,
        text width=5.6cm] at (0.1,-1.25)
       {\textbf{Three containers on disk:} 139 MB shared, plus whatever each
        one has written. Three virtual machines: three full root
        filesystems.};
\end{tikzpicture}}%
{Overlayfs. The layers are why pulling a second image from the same base
downloads almost nothing, and why a thousand containers of one image cost
little more on disk than one.}{overlayfs}

\begin{worked}[Density and Start-Up, Costed]
A host of 64 cores and 512 GB must run as many copies as possible of one
service that needs 256 MB of memory. Compare virtual machines with containers.

\smallskip
\textbf{Step 1 --- memory per instance, as a virtual machine.}
\begin{tight}
  \item the application: 256 MB;
  \item a guest Linux that can run it: 512 MB is a realistic floor;
  \item per-guest hypervisor overhead --- the VMCS, the EPT tables, the device
        models of \chref{paravirt}: call it 100 MB.
\end{tight}
\[ 256 + 512 + 100 = 868 \text{ MB per guest} \]

\smallskip
\textbf{Step 2 --- memory per instance, as a container.}
\begin{tight}
  \item the application: 256 MB;
  \item no guest kernel --- the host's is shared;
  \item runtime overhead: a few MB. Call it 4 MB.
\end{tight}
\[ 256 + 4 = 260 \text{ MB per container} \]

\smallskip
\textbf{Step 3 --- how many fit.} Reserve 32 GB for the host, leaving 480 GB:
\[ \text{VMs} = \frac{480 \times 1024}{868} = 566 \qquad
   \text{containers} = \frac{480 \times 1024}{260} = \mathbf{1{,}890} \]
A ratio of \textbf{3.3 : 1} --- and note it is not the hundredfold figure
containers are often credited with, because the \emph{application's} 256 MB
dominates both sides. The container advantage is large exactly when the
workload is small.

\smallskip
\textbf{Step 4 --- the disk, where the ratio is larger.} 566 guests each with a
2 GB root disk is 1.1 TB, even thin-provisioned. 1{,}890 containers from one
139 MB image share that image: 139 MB plus each container's writes. Call it
10 MB each:
\[ 139 + 1{,}890 \times 10 = 19 \text{ GB, against } 1{,}100 \text{ GB}
   \;\approx\; \mathbf{58 : 1} \]

\smallskip
\textbf{Step 5 --- start-up, where it is larger still.} A guest boots a kernel,
initialises virtual hardware and starts an init system: 20--40 s. A container
is \texttt{clone()} with some flags, a mount and an \texttt{execve}: 50--200 ms.
\[ \frac{30 \text{ s}}{0.1 \text{ s}} = \mathbf{300 : 1} \]

\smallskip
\textbf{Step 6 --- what the three ratios are actually telling you.} Memory
density, 3:1, is the least impressive number and the one usually quoted. The
ones that changed how software is deployed are disk and start-up: when an
instance costs 10 MB and starts in a tenth of a second, you can afford to run
one per request, scale on demand, and redeploy continuously. \chref{kubernetes}
is what people built once that became true.
\end{worked}

\begin{alertbox}[What the ratios cost: one kernel, one failure]
Every container on that host shares one kernel. A kernel panic takes all 1{,}890
down; a kernel privilege-escalation bug gives any one of them the host. The
Linux system-call interface is around 350 calls wide and every one is reachable
from inside a container, where a virtual machine's guest would have to escape
its own kernel first and then the hypervisor.

\smallskip
That is the trade, stated numerically: 3$\times$ the density, 300$\times$ the
start-up speed, and an isolation boundary enforced by software rather than by
the processor. \chref{lightweight} is the chapter where people decline to
choose.
\end{alertbox}

\begin{pitfall}
\begin{tight}
  \item \textbf{Calling a container a lightweight virtual machine.} There is no
        virtual machine and no guest kernel. The phrase predicts the wrong
        things about security, about kernel modules, and about what
        \texttt{uname} reports inside.
  \item \textbf{Expecting a container to run a different kernel.} It cannot. A
        container ``running CentOS'' on an Ubuntu host is CentOS \emph{user
        space} on Ubuntu's kernel. Nothing needing a specific kernel version or
        a kernel module will work.
  \item \textbf{Running an application directly as PID 1 and wondering about
        zombies or ten-second stops.} Use an init, and know why.
  \item \textbf{Setting CPU limits carefully and memory limits by guesswork.}
        CPU throttles; memory kills. Exit code 137 is the OOM killer and it is
        the commonest container incident there is.
  \item \textbf{Writing large amounts of data to the container's writable
        layer.} Copy-up through overlayfs is slow, and the layer disappears
        with the container. Mount a volume.
  \item \textbf{Assuming a container is isolated because it has a private
        namespace.} Namespaces hide things; they do not reduce the kernel
        attack surface. \chref{security} separates the two properly.
\end{tight}
\end{pitfall}

\begin{lab}[Building a Container by Hand]
No Docker. Builds an isolated process tree from the kernel primitives, so that
the next chapter's tooling is recognisable rather than magical. Needs a Linux
host and root. Expect thirty minutes.

\begin{lstlisting}[language=bash]
#!/usr/bin/env bash
# ch08_byhand.sh -- namespaces, cgroups and overlayfs, without a runtime.
set -euo pipefail
ROOT=/tmp/vsroot

# --- 1. What namespaces does an ordinary process already have? ---------
# Every process is in a namespace; a container just gets different ones.
ls -l /proc/self/ns/

# --- 2. A root filesystem to be the container's world. -----------------
sudo apt-get install -y debootstrap util-linux
sudo mkdir -p "$ROOT"
sudo debootstrap --variant=minbase jammy "$ROOT" \
     http://archive.ubuntu.com/ubuntu/

# --- 3. Enter new namespaces. This is the whole mechanism. -------------
# --pid --mount --uts --ipc --net: one flag per namespace from the table.
sudo unshare --pid --mount --uts --ipc --net --fork --mount-proc \
     chroot "$ROOT" /bin/bash -c '
       hostname container-1
       mount -t proc proc /proc
       echo "--- inside ---"
       hostname            # the UTS namespace
       ps aux              # the PID namespace: we are PID 1, host invisible
       ip addr             # the network namespace: only lo, and it is down
       exit'

# --- 4. Prove the host is unaffected. ----------------------------------
hostname ; ps aux | wc -l

# --- 5. Limit it with a control group. ---------------------------------
# cgroup v2: one unified tree under /sys/fs/cgroup.
sudo mkdir -p /sys/fs/cgroup/vsdemo
echo "+cpu +memory" | sudo tee /sys/fs/cgroup/cgroup.subtree_control
echo "50000 100000" | sudo tee /sys/fs/cgroup/vsdemo/cpu.max   # half a core
echo "100M"         | sudo tee /sys/fs/cgroup/vsdemo/memory.max

# Put a busy process in it and watch it be throttled, not killed.
sudo bash -c 'echo $$ > /sys/fs/cgroup/vsdemo/cgroup.procs; \
              timeout 10 yes > /dev/null' &
sleep 11
cat /sys/fs/cgroup/vsdemo/cpu.stat | grep throttled

# --- 6. Now cross the memory limit, and watch it be killed. -----------
# CPU throttles; memory kills. This is the asymmetry of section 8.3.
sudo bash -c 'echo $$ > /sys/fs/cgroup/vsdemo/cgroup.procs; \
              head -c 200M /dev/zero | tail -c 1' || \
  echo "killed, exit $? -- 137 is 128+9, the OOM killer"
dmesg | tail -3 | grep -i 'out of memory' || true

# --- 7. Share one root between two containers with overlayfs. ---------
sudo mkdir -p /tmp/ovl/{upper1,upper2,work1,work2,m1,m2}
for i in 1 2; do
  sudo mount -t overlay overlay \
    -o lowerdir="$ROOT",upperdir=/tmp/ovl/upper$i,workdir=/tmp/ovl/work$i \
    /tmp/ovl/m$i
done
# Write in one; the other and the shared lower layer are untouched.
sudo touch /tmp/ovl/m1/etc/only-in-one
ls /tmp/ovl/m1/etc/only-in-one
ls /tmp/ovl/m2/etc/only-in-one 2>&1 | tail -1
sudo du -sh /tmp/ovl/upper1 /tmp/ovl/upper2 "$ROOT"

# --- 8. Tear down. ----------------------------------------------------
for i in 1 2; do sudo umount /tmp/ovl/m$i || true; done
sudo rmdir /sys/fs/cgroup/vsdemo || true
sudo rm -rf /tmp/ovl "$ROOT"
\end{lstlisting}

\textbf{What to notice.} Step 3 is a container. There is no daemon, no image
format and no runtime --- one command, five flags and a directory. Everything
\chref{docker} adds is packaging and distribution on top of exactly this.
\end{lab}

\begin{checkpoint}
\begin{tightnum}
  \item A container restarts every few hours with exit code 137. Say what that
        code means, which limit is responsible, and why the symptom is a kill
        rather than a slowdown.
  \item A team wants to run a container that loads a kernel module for a custom
        device. Explain why this cannot work as stated, and give the two
        options available to them.
  \item Three containers run from the same 400 MB image, and each writes 50 MB
        of logs into its own filesystem. How much disk is used, and what single
        change would stop the logs counting at all?
\end{tightnum}
\end{checkpoint}

\begin{chaptersummary}
\begin{tight}
  \item A container cuts the stack above the kernel instead of below it: one
        kernel, many private views of it. Everything else follows from that.
  \item Eight namespaces isolate what a process can see --- mount, UTS, IPC,
        PID, network, user, cgroup, time. The user namespace is what makes
        rootless containers possible.
  \item PID 1 is special: it must reap orphans and ignores default-action
        signals. Run an init unless you have handled both.
  \item Control groups bound what a process tree may use. CPU limits throttle;
        memory limits kill. The asymmetry is the same one as in \chref{cpumem}.
  \item Overlayfs stacks read-only layers under one writable layer, with
        copy-up on write, so many containers share one image.
  \item Measured against virtual machines: roughly 3:1 on memory, 58:1 on disk,
        300:1 on start-up. The last two are what changed how software is
        deployed.
  \item The cost is a single shared kernel and a 350-call attack surface.
\end{tight}
\end{chaptersummary}

\begin{reviewq}
  \item \textbf{(MCQ)} The namespace that gives a container its own filesystem
        tree is: (a)~PID; (b)~mount; (c)~UTS; (d)~user.
  \item \textbf{(MCQ)} A container exceeding its memory limit is: (a)~throttled;
        (b)~swapped by the host; (c)~has a process killed by the OOM killer;
        (d)~given memory from another container.
  \item \textbf{(Short)} Explain copy-up, and say what it implies about writing
        large files inside a container.
  \item \textbf{(Short)} Why can a container not run a different kernel from
        its host, and what is actually different about a ``CentOS container''
        on an Ubuntu host?
  \item \textbf{(Short)} State the two behaviours that make PID 1 special, and
        the symptom each produces when ignored.
  \item \textbf{(Applied)} Repeat the worked example for a service needing
        2 GB rather than 256 MB, on the same host. Give the three ratios and
        explain why the memory ratio changed so much while the start-up ratio
        did not.
  \item \textbf{(Applied)} A team proposes moving forty internal services from
        virtual machines to containers on shared hosts. Write the three
        questions you would ask before agreeing, and say what answer to each
        would make you recommend \chref{lightweight}'s approach instead.
\end{reviewq}
